\documentclass[11pt, oneside]{article}   	% use "amsart" instead of "article" for AMSLaTeX format
\usepackage{geometry}                		% See geometry.pdf to learn the layout options. There are lots.
\geometry{letterpaper}                   		% ... or a4paper or a5paper or ... 
%\geometry{landscape}                		% Activate for rotated page geometry
%\usepackage[parfill]{parskip}    		% Activate to begin paragraphs with an empty line rather than an indent
\usepackage{graphicx}				% Use pdf, png, jpg, or eps§ with pdflatex; use eps in DVI mode
								% TeX will automatically convert eps --> pdf in pdflatex		
\usepackage{amssymb}
\usepackage{amsmath}
\usepackage{url}

%SetFonts

%SetFonts


\title{Equations to be implemented in Castro}
\author{Remi Lehe}
%\date{}							% Activate to display a given date or no date

\begin{document}
\maketitle
%\section{}
%\subsection{}

\section{Derivation}

We start from the equations in Mewes at al. (\url{https://arxiv.org/abs/2305.16779}) and derive the corresponding form that is compatible with the Castro set of equations (\url{https://amrex-astro.github.io/Castro/docs/Hydrodynamics.html#conservation-forms}).

\subsection{Conservation of mass}

Equation (1) from Mewes et al. is already in the right form:
\[ \partial_t \rho + \boldsymbol{\nabla}\cdot( \rho\boldsymbol{v} ) = 0  \]
This implies that the corresponding source term in Castro is 0:
\[ S_{ext,\rho} = 0 \] 

\subsection{Evolution of the population of each species}

We have 3 species: $H$, $H^+$ and $e^{-}$. Since the plasma is assumed to be neutral ($n_e = n_{H^+}$), we only need to track the mass fraction of $H$ and $H^+$ ($X_H$ and $X_{H^+}$), and $n_e$ is then given by
\[ n_e = \rho X_{H^+} / m_H \]
where $m_H$ is the mass of one ion (or atom) of $H$.

To find the evolution of the fractions $X_k$, we start from the left-hand side of the Castro form and use differentiating rules for $\rho \times X_k$:
\[ \partial_t \rho X_k + \boldsymbol{\nabla}\cdot( \rho X_k \boldsymbol{v} ) =  (\rho\partial_t + \rho \boldsymbol{v}\cdot\boldsymbol{\nabla})X_k + (\partial_t \rho + \boldsymbol{\nabla}\cdot(\rho \boldsymbol{v} ) ) X_k\]
The first term on the right-hand side corresponds to Eq. (5) from Mewes et al. (with $w_\alpha = X_k$) and the second term on the right-hand side is 0 (due to the conservation of mass). \textbf{We furthermore neglect diffusion in Eq. (5) from Mewes et al.}
\[ \partial_t \rho X_k + \boldsymbol{\nabla}\cdot( \rho X_k \boldsymbol{v} ) = R_k \]

More specifically, $R_k$ is obtained from appendix D in Mewes et al:
\textbf{One needs to be careful regarding units, here. Mewes et al. use $m^3.s^{-1}$ for $A_j$ whereas Castro uses CGS.}
\begin{align}
R_{H+} = S_{ext, \rho X_{H+} } &= m_H \sum_j A_j \left( \frac{k_BT_e}{e} \right)^{q_j}\exp\left( - \frac{\epsilon_j}{k_B T_e} \right)\left[ n_H n_e - n_{H^+}n_e^2 \lambda^3 e^{\epsilon_{ion}/k_BT_e} \right] \\
&= \frac{\rho^2}{m_H}\sum_j A_j \left( \frac{k_BT_e}{e} \right)^{q_j}\exp\left( - \frac{\epsilon_j}{k_B T_e} \right)\left[ X_H X_{H^+} - X_{H^+}^3 \frac{\rho \lambda^3}{m_H} e^{\epsilon_{ion}/k_BT_e} \right] 
\end{align}
\[ R_H = S_{ext, \rho X_{H} } = -R_{H+} \]
where: \[ \lambda = \sqrt{\frac{h^2}{2\pi m_e k_B T_e} } \]

\subsection{Conservation of momentum}

We start from the left-hand side of the Castro form, and use differentiating rules for $\rho \times \boldsymbol{v}$:
\[
\partial_t \rho \boldsymbol{v} + \boldsymbol{\nabla}\cdot(\rho \boldsymbol{v} \boldsymbol{v}) = (\rho\partial_t + \rho \boldsymbol{v}\cdot\boldsymbol{\nabla}) \boldsymbol{v} + (\partial_t \rho + \boldsymbol{\nabla}\cdot(\rho \boldsymbol{v} ) ) \boldsymbol{v}
\]
The first term on the right-hand side corresponds to Eq. (2) from Mewes et al. and the second term on the right-hand side is 0 (due to the conservation of mass). \textbf{We furthermore neglect the viscosity and external forces in Eq. (2) from Mewes et al.}
\[ \partial_t \rho \boldsymbol{v} + \boldsymbol{\nabla}\cdot(\rho \boldsymbol{v} \boldsymbol{v}) = -\boldsymbol{\nabla}p \]
This implies that the corresponding source term in Castro is again 0:
\[ \boldsymbol{S}_{ext,\rho\boldsymbol{u}} = \boldsymbol{0} \]

\subsection{Conservation of energy}

We have a separate equation for the heavy species and the electrons. For the heavy species, i.e. $H$ and $H^+$ together (Eq. 3 in Mewes et al.):
\[ \rho C_h (\partial_t + \boldsymbol{v}\cdot\boldsymbol{\nabla})T_h = \boldsymbol{\nabla}\cdot(\lambda_h\boldsymbol{\nabla}T_h) + (\partial_t + \boldsymbol{v}\cdot\boldsymbol{\nabla})p_h + n_e\nu_{eh}^\epsilon \frac{3}{2}k_b(T_e - T_h) \]
where $C_h = \frac{5}{2}\frac{k_B}{m_H}$ and $p_H = \frac{\rho k_B T_h}{m_H}$.

Since $C_h$ is constant, we have:
\begin{align*} 
\rho C_h (\partial_t + \boldsymbol{v}\cdot\boldsymbol{\nabla})T_h &=  (\partial_t + \boldsymbol{v}\cdot\boldsymbol{\nabla}) \rho C_h T_h - C_h T_h  (\partial_t + \boldsymbol{v}\cdot\boldsymbol{\nabla})\rho \\
&= (\partial_t \rho C_h T_h + \boldsymbol{\nabla}\cdot(\rho C_h T_h \boldsymbol{v}) - \rho C_h T_h \boldsymbol{\nabla}\cdot\boldsymbol{v} ) + C_hT_h \rho \boldsymbol{\nabla}\cdot\boldsymbol{v} \\ 
& =  \partial_t \rho C_h T_h + \boldsymbol{\nabla}\cdot(\rho C_h T_h \boldsymbol{v})
\end{align*}

We can thus rewrite the equation for heavy species as:
\[ \partial_t \rho C_h T_h + \boldsymbol{\nabla}\cdot(\rho C_h T_h \boldsymbol{v}) = \boldsymbol{\nabla}\cdot(\lambda_h\boldsymbol{\nabla}T_h) + (\partial_t p_h + \boldsymbol{\nabla}\cdot(p_h \boldsymbol{v}) - p_h \boldsymbol{\nabla}\cdot \boldsymbol{v}) + n_e\nu_{eh}^\epsilon \frac{3}{2}k_b(T_e - T_h) \]

And finally, using
\[ e_h = C_h T_e - \frac{p_h}{\rho} = \frac{3}{2}\frac{k_B T_h}{m_H} \]
we have:
\[ \partial_t \rho e_h + \boldsymbol{\nabla}\cdot(\rho e_h \boldsymbol{v}) = \boldsymbol{\nabla}\cdot(\lambda_h\boldsymbol{\nabla}T_h) - p_h \boldsymbol{\nabla}\cdot \boldsymbol{v} + n_e\nu_{eh}^\epsilon \frac{3}{2}k_b(T_e - T_h) \]


For the electron species:
\[ \rho C_e (\partial_t + \boldsymbol{v}\cdot\boldsymbol{\nabla})T_e = \boldsymbol{\nabla}\cdot(\lambda_e\boldsymbol{\nabla}T_e) + (\partial_t + \boldsymbol{v}\cdot\boldsymbol{\nabla})p_e - n_e\nu_{eh}^\epsilon \frac{3}{2}k_b(T_e - T_h) + Q_e\]
where $C_e = \frac{5}{2}\frac{X_{H^+}k_B}{m_H}$ and $p_e = \frac{\rho X_{H^+} k_B T_e}{m_H}$.

Let us re-express the left-hand side:
\begin{align*} 
\rho \frac{5}{2}\frac{X_{H^+}k_B}{m_H} (\partial_t + \boldsymbol{v}\cdot\boldsymbol{\nabla})T_e  &=  (\partial_t + \boldsymbol{v}\cdot\boldsymbol{\nabla}) \rho \frac{5}{2}\frac{X_{H^+}k_B T_e}{m_H} - \frac{5}{2}\frac{k_BT_e}{m_H} (\partial_t + \boldsymbol{v}\cdot\boldsymbol{\nabla}) \rho X_{H^+} \\
&= \partial_t \rho C_e T_e + \boldsymbol{\nabla}\cdot( \rho C_e T_e \boldsymbol{v} ) - \rho C_e T_e\boldsymbol{\nabla}\cdot \boldsymbol{v}  - \frac{5}{2}\frac{k_BT_e}{m_H} \left[ R_{H^+} - \rho X_{H^+}\boldsymbol{\nabla}\cdot\boldsymbol{v}\right] \\
& = \partial_t \rho C_e T_e + \boldsymbol{\nabla}\cdot( \rho C_e T_e \boldsymbol{v} ) - Q_e^{num}
\end{align*}

We can thus rewrite the equation for the electrons as:
\[ \partial_t \rho C_e T_e + \boldsymbol{\nabla}\cdot(\rho C_e T_e \boldsymbol{v}) = \boldsymbol{\nabla}\cdot(\lambda_h\boldsymbol{\nabla}T_e) + (\partial_t + \boldsymbol{v}\cdot\boldsymbol{\nabla})p_e - n_e\nu_{eh}^\epsilon \frac{3}{2}k_b(T_e - T_h) + Q_e^{chem} \]

And finally, using
\[ e_h = C_e T_e - \frac{p_e}{\rho} = \frac{3}{2}X_{H^+}\frac{k_B T_e}{m_H} \]
we have:
\[ \partial_t \rho e_e + \boldsymbol{\nabla}\cdot(\rho e_e \boldsymbol{v}) = \boldsymbol{\nabla}\cdot(\lambda_e\boldsymbol{\nabla}T_e) - p_e\boldsymbol{\nabla}\cdot\boldsymbol{v} - n_e\nu_{eh}^\epsilon \frac{3}{2}k_b(T_e - T_h) + Q_e^{chem} \]

Finally, by summing the equations for the electrons and heavy species, we have:
\[  \partial_t \rho e + \boldsymbol{\nabla}\cdot(\rho e \boldsymbol{v}) = \boldsymbol{\nabla}\cdot(\lambda_h\boldsymbol{\nabla}T_e) + \boldsymbol{\nabla}\cdot(\lambda_h\boldsymbol{\nabla}T_h) - p\boldsymbol{\nabla}\cdot\boldsymbol{v} + Q_e^{chem} \]
with:
\[ e = e_h + e_e =  \frac{3}{2}\frac{k_B}{m_H}(T_h + X_{H^+}T_e) \qquad p = p_h + p_e = \frac{\rho k_B}{m_H}(T_h + X_{H^+}T_e) \]

\textbf{If we neglect thermal diffusion}, this corresponds to \url{https://amrex-astro.github.io/Castro/docs/Hydrodynamics.html#internal-energy-and-temperature} with:
\[ S_{ext, E} = Q_e^{chem} \]

\section{Final set of equations}

When neglecting diffusion, viscosity and thermal diffusion:

\begin{align}
\partial_t \rho + \boldsymbol{\nabla}\cdot( \rho\boldsymbol{v} ) &= 0 &\qquad S_{ext,\rho} = 0 \\
\partial_t \rho \boldsymbol{v} + \boldsymbol{\nabla}\cdot(\rho \boldsymbol{v} \boldsymbol{v}) &= -\boldsymbol{\nabla}p &\qquad \boldsymbol{S}_{ext,\rho\boldsymbol{u}} = \boldsymbol{0} \\
 \partial_t \rho E + \boldsymbol{\nabla}\cdot(\rho E \boldsymbol{v}) &= - \boldsymbol{\nabla}\cdot (p\boldsymbol{v}) -R_{H+} \frac{\epsilon_{ion}}{m_H}& \qquad S_{ext, E} = -R_{H+} \frac{\epsilon_{ion}}{m_H} \\
 \partial_t \rho X_H + \boldsymbol{\nabla}\cdot( \rho X_H \boldsymbol{v} ) &= -R_{H+} & \qquad S_{ext, X_H} = - R_{H^+}\\
 \partial_t \rho X_{H^+} + \boldsymbol{\nabla}\cdot( \rho X_{H^+} \boldsymbol{v} ) &= R_{H+} & \qquad S_{ext, X_{H^+}} = R_{H^+}\\
  \partial_t \rho e_h + \boldsymbol{\nabla}\cdot(\rho e_h \boldsymbol{v}) &= - p_h \boldsymbol{\nabla}\cdot \boldsymbol{v} + \rho X_{H^+}\nu_{eh}^\epsilon \frac{3}{2}\frac{k_B}{m_e}(T_e - T_h) 
\end{align}

with 
\[ R_{H+} = \frac{\rho^2}{m_H}\sum_j A_j \left( \frac{k_BT_e}{e} \right)^{q_j}\exp\left( - \frac{\epsilon_j}{k_B T_e} \right)\left[ X_H X_{H^+} - X_{H^+}^3 \frac{\rho \lambda^3}{m_H} e^{\epsilon_{ion}/k_BT_e} \right]  \]
and
\[ e = e_h + e_e =  \frac{3}{2}\frac{k_B}{m_H}(T_h + X_{H^+}T_e) \qquad p = p_h + p_e = \frac{\rho k_B}{m_H}(T_h + X_{H^+}T_e) \]

\end{document}  